\section{Теоретическое обоснование мезоскопической модели транспортного потока}

Отказ от микроскопического моделирования каждой отдельной машины ставит вопрос о математическом аппарате, который сможет достоверно описывать физику заторов. В основе разрабатываемого решения лежат принципы классической гидродинамической теории транспортного потока, переложенные на архитектуру высокопроизводительного графового движка.

\subsection{Эйлеров подход и модель пространственных очередей}

Классические микроскопические симуляторы (подобные SUMO) используют \textit{лагранжев подход}, отслеживая непрерывные координаты $x(t), y(t)$ каждого транспортного средства. При росте количества агентов $N$ сложность выявления коллизий (потенциальных столкновений на перекрестках) возрастает как $O(N^2)$, что делает симуляцию мегаполиса практически невыполнимой в реальном времени.

В разработанной системе применен \textbf{эйлеров подход}. Физическое пространство рассматривается не как непрерывная плоскость, а как набор неподвижных гранулированных сечений (ребер графа). Внимание симулятора концентрируется не на отдельном автомобиле, а на макроскопическом параметре — скалярной плотности потока, проходящего через это сечение.

Движение описывается моделью пространственных очередей (Queue-based kinematics). В любой момент времени $t$ каждое направленное ребро $e$ характеризуется тремя фундаментальными константами:
\begin{itemize}
    \item $C_{max}$ (Jam Capacity) — максимальная пространственная вместимость ребра (штук ТС);
    \item $V(t)$ (Volume) — текущее количество находящихся на сегменте агентов;
    \item $T_{free}$ (Free Flow Time) — идеальное время проезда сегмента при нулевом трафике (определяемое скоростным лимитом и длиной ребра).
\end{itemize}

Агенты перемещаются из очереди одного ребра в очередь следующего ребра в соответствии с их маршрутом, при условии, что целевое ребро еще имеет свободное пространство ($V(t) < C_{max}$).

\subsection{Теория Spillback и гидродинамическое приближение LWR}

Ключевой особенностью реальных заторов является их способность распространяться против движения транспортного потока. Этот эффект, известный как \textit{Spillback}, возникает, когда улица полностью заполняется автомобилями, блокируя въезд на нее с пересекающих магистралей, что, в свою очередь, парализует и их.

Для моделирования физики Spillback применяется дискретная аппроксимация гидродинамической теории транспортного потока Лайтхилла — Уизема — Ричардса (Lighthill-Whitham-Richards, LWR). 
В контексте эйлерова подхода модель LWR реализуется через жесткий барьер вместимости: если количество транспортных средств на ребре достигает его критической плотности ($V(t) \ge C_{max}$), пропускная способность входного сечения ребра (Inflow Capacity) мгновенно падает до нуля. 

Любой агент, попытавшийся въехать на переполненное ребро, останавливается и остается на предыдущем ребре графа. Тем самым агент занимает место на предыдущем ребре, увеличивая его плотность $V_{prev}(t)$. Как только $V_{prev}(t)$ достигнет своего лимита $C^{prev}_{max}$, волна затора перекинется еще на один уровень назад. Данная физика позволяет формировать достоверные «древовидные» заторы на картах.

Важной оптимизацией, примененной в системе, является выполнение всех расчетов плотности Spillback строго в пространстве агентов (\textbf{Agent Space}), без предварительного масштабирования через коэффициент симуляции (Agent Scale Factor). Применение Agent Scale Factor на этапе вычисления емкости ребер ($C_{max} / ASF$) приводило к катастрофическим ошибкам округления (целочисленному квантованию) на коротких ребрах. Отказ от этого масштабирования при расчете Spillback-блокировок предотвратил накопление численных дедлоков и снизил долю перманентно заблокированных зон («черных зон») в симулируемой сети с 22,3\% до 5,7\%.

\subsection{Динамические функции сопротивления сети (модель BPR)}

В условиях плотного, но еще не заблокированного потока ($0 < V(t) < C_{max}$), скорость движения автомобилей неизбежно снижается из-за взаимных помех. Связь между плотностью потока и временем задержки описывается нелинейной функцией сопротивления. В разрабатываемой системе применена модифицированная функция Бюро дорог общего пользования США (BPR — Bureau of Public Roads).

Полное время проезда сегмента $W_{total}(t)$ в любой момент времени $t$ вычисляется по формуле:
\begin{equation}
    W_{total}(t) = T_{free} + \Delta W_{dynamic}(t) + mpr\_penalty
\end{equation}
где $T_{free}$ — базовая нижняя граница времени проезда (учитывающая скоростной лимит и константные штрафы за сложные геометрические маневры, например левые повороты на перекрестках), $mpr\_penalty$ — искусственный директивный штраф, выставляемый Модулем принятия решений для отвода потока от зоны инцидента, а $\Delta W_{dynamic}(t)$ — динамическая задержка затора.

В классической теории функция BPR имеет вид многочлена 4-й степени, однако для городских сетей сверхбольшой размерности вычисление плавающей запятой в горячем цикле роутера недопустимо. Функция была преобразована к полиному 2-й степени ($\beta = 2$) с адаптивным коэффициентом крутизны $\alpha$:
\begin{equation}
    \Delta W_{dynamic}(t) = T_{free} \cdot \alpha \cdot \left(\frac{V(t)}{C_{300}}\right)^2
\end{equation}

Для обеспечения максимальной утилизации конвейера процессора и полного исключения дорогостоящих операций деления в горячем цикле алгоритма A*, все постоянные коэффициенты ребра сворачиваются в предрассчитанный 32-битный множитель $K_{magic}$ с использованием арифметики с фиксированной точкой (битовый сдвиг на 20 разрядов влево):
\begin{equation}
    K_{magic} = \left\lfloor \frac{T_{free} \cdot \left( \frac{V_{free}}{5} - 1 \right)}{C_{300}^2} \cdot 2^{20} \right\rfloor
\end{equation}
где $C_{300}$ — емкость ребра для интервала 5 минут (300 секунд), а коэффициент $\alpha = \frac{V_{free}}{5} - 1$ аналитически компенсирует глобальное снижение степени BPR-функции, сохраняя достоверный профиль падения скорости.

Таким образом, в критическом по производительности вычислительном цикле роутера расчет сложной нелинейной BPR-задержки сводится всего к трем тактам процессора: одному целочисленному умножению, возведению в квадрат и быстрому битовому сдвигу вправо:
\begin{equation}
    \Delta W_{dynamic}(t) = (K_{magic} \cdot V(t)^2) \gg 20
\end{equation}
